
\end{equation}
où @@M0@@. Par l'hypothèse de récurrence on a @@M1@@, d'où, compte tenu des formules (1.3.8) et (1.3.9), l'assertion
@@M2@@

\textbf{Remarque 1.4.} Soit @@M3@@ un schéma localement noethérien. On notera que nous avons fait usage des \emph{seules} propriétés suivantes de la « fonction dimension » @@M4@@ sur les @@M5@@-schémas propres @@M6@@ au cours de cette section :
\begin{enumerate}[label=(\roman*)]
\item Pour tout @@M7@@-schéma propre @@M8@@ et pour tout @@M9@@, on a @@M10@@.
\item Si @@M11@@ est une spécialisation de @@M12@@, on a
@@M13@@
\item Pour tout @@M14@@-morphisme @@M15@@ de @@M16@@-schémas propres et pour tout @@M17@@, on a
@@M18@@
\item Il y a égalité dans (iii), si @@M19@@ est une immersion ou si @@M20@@ est irréductible, @@M21@@ birationnel et surjectif, et @@M22@@ le point générique de @@M23@@.
\end{enumerate}

\uline{Cela signifie que tous les résultats de cette section, en particulier} (1.3.2) \uline{et} (1.3.3)\uline{, restent valables tels quels, si l'on y remplace la} « \uline{dimension} » \uline{par une fonction} @@M24@@ \uline{satisfaisant aux conditions} (i)à(iv) \uline{ci-dessus et si l'on définit la filtration de} @@M25@@ \uline{suivant une telle fonction} @@M26@@.

